% =============================================================================
% DIAPOSITIVAS — PROYECTO 4: MUESTRA INTELIGENTE
% Estilo: minimalista, juvenil, poco texto
% Compilar en Overleaf con pdfLaTeX
% =============================================================================

\documentclass[aspectratio=169, 10pt]{beamer}

% ── Tema base ─────────────────────────────────────────────────────────────────
\usetheme{metropolis}
\usepackage{appendixnumberbeamer}

% ── Paquetes ──────────────────────────────────────────────────────────────────
\usepackage[spanish]{babel}
\usepackage[utf8]{inputenc}
\usepackage[T1]{fontenc}
\usepackage{booktabs}
\usepackage{colortbl}
\usepackage{xcolor}
\usepackage{tcolorbox}
\usepackage{fontawesome5}
\usepackage{listings}
\usepackage{tikz}
\usepackage{graphicx}
\usetikzlibrary{positioning}

% ── Colores ───────────────────────────────────────────────────────────────────
\definecolor{cdark}{HTML}{111111}
\definecolor{cverde}{HTML}{006847}
\definecolor{camarillo}{HTML}{FDD000}
\definecolor{cgris}{HTML}{999999}
\definecolor{ccard}{HTML}{1E1E1E}
\definecolor{clight}{HTML}{F5F5F5}
\definecolor{ccode}{HTML}{161616}

% ── Colores del tema ──────────────────────────────────────────────────────────
\setbeamercolor{normal text}{fg=white, bg=cdark}
\setbeamercolor{frametitle}{bg=cdark, fg=white}
\setbeamercolor{progress bar}{fg=camarillo, bg=cdark}
\setbeamercolor{title separator}{fg=camarillo}
\setbeamercolor{alerted text}{fg=camarillo}
\setbeamercolor{block title}{bg=cverde, fg=white}
\setbeamercolor{block body}{bg=ccard, fg=white}
\setbeamercolor{background canvas}{bg=cdark}

\metroset{block=fill, progressbar=foot, titleformat=smallcaps}

% ── Sin logo ni pie redundante ────────────────────────────────────────────────
\setbeamertemplate{frame numbering}[fraction]

% ── tcolorbox ─────────────────────────────────────────────────────────────────
\tcbuselibrary{skins}
\newtcolorbox{tarjeta}{
  colback=ccard, colframe=white!10!ccard,
  arc=8pt, boxrule=0.5pt,
  left=8pt, right=8pt, top=6pt, bottom=6pt,
}
\newtcolorbox{tarjetaverde}{
  colback=cverde!20!cdark, colframe=cverde,
  arc=8pt, boxrule=0.8pt,
  left=8pt, right=8pt, top=6pt, bottom=6pt,
}
\newtcolorbox{formula}{
  colback=ccode, colframe=camarillo,
  arc=6pt, boxrule=0pt,
  borderline west={3pt}{0pt}{camarillo},
  left=10pt, right=10pt, top=6pt, bottom=6pt,
}

% ── Código Python ─────────────────────────────────────────────────────────────
\lstdefinestyle{py}{
  language=Python,
  basicstyle=\ttfamily\fontsize{6.8}{9}\selectfont\color{white},
  keywordstyle=\color{camarillo}\bfseries,
  commentstyle=\color{cgris}\itshape,
  stringstyle=\color{green!60!white},
  backgroundcolor=\color{ccode},
  frame=none,
  showstringspaces=false,
  breaklines=true,
  tabsize=2,
}

% ── Metadatos ─────────────────────────────────────────────────────────────────
\title{\textbf{Muestra Inteligente}}
\subtitle{Backtracking aplicado a selección de muestras representativas}
\author{Sebastián Gutiérrez \and Mateo Bernal \and Thomas Ramírez}
\institute{Universidad Nacional de Colombia — Sede Manizales\\
           Introducción a las Ciencias de la Computación\\
           Prof.\ Carlos Manuel Orrego Franco}
\date{17 de junio de 2026}

% =============================================================================
\begin{document}
% =============================================================================

% ── PORTADA ───────────────────────────────────────────────────────────────────
\begin{frame}
  \titlepage
\end{frame}


% ── 1. EL PROBLEMA ────────────────────────────────────────────────────────────
{
\setbeamercolor{background canvas}{bg=clight}
\setbeamercolor{normal text}{fg=cdark}
\setbeamercolor{frametitle}{bg=clight, fg=cdark}
\begin{frame}{El problema}
  \begin{center}
    {\large ¿Cómo elegir \textcolor{cverde}{\textbf{6 estudiantes}} sin sesgar la muestra?}\\[0.5cm]
    {\normalsize Bienestar Universitario necesita una muestra \textbf{pequeña y representativa}.}
    \vspace{0.6cm}

    \begin{columns}[c]
      \begin{column}{0.3\textwidth}
        \begin{center}
          {\fontsize{36}{40}\selectfont\bfseries\color{cverde} 20}\\
          {\small\color{cgris} estudiantes en el pool}
        \end{center}
      \end{column}
      \begin{column}{0.3\textwidth}
        \begin{center}
          {\fontsize{36}{40}\selectfont\bfseries\color{cverde} 6}\\
          {\small\color{cgris} a seleccionar}
        \end{center}
      \end{column}
      \begin{column}{0.3\textwidth}
        \begin{center}
          {\fontsize{36}{40}\selectfont\bfseries\color{cverde} 38K}\\
          {\small\color{cgris} combinaciones posibles}
        \end{center}
      \end{column}
    \end{columns}
  \end{center}
\end{frame}
}


% ── 2. POR QUÉ BACKTRACKING ───────────────────────────────────────────────────
\begin{frame}{Por qué \alert{backtracking}}
  \vspace{0.3cm}
  \begin{columns}[c]
    \begin{column}{0.3\textwidth}
      \begin{tarjeta}
        \begin{center}
          {\Large\faHammer\par}
          \vspace{4pt}
          {\small\bfseries Construye}\\[2pt]
          {\footnotesize\color{cgris} la muestra paso a paso}
        \end{center}
      \end{tarjeta}
    \end{column}
    \begin{column}{0.3\textwidth}
      \begin{tarjeta}
        \begin{center}
          {\Large\faCut\par}
          \vspace{4pt}
          {\small\bfseries Poda}\\[2pt]
          {\footnotesize\color{cgris} ramas imposibles antes de explorarlas}
        \end{center}
      \end{tarjeta}
    \end{column}
    \begin{column}{0.3\textwidth}
      \begin{tarjeta}
        \begin{center}
          {\Large\faUndo\par}
          \vspace{4pt}
          {\small\bfseries Retrocede}\\[2pt]
          {\footnotesize\color{cgris} cuando viola una cuota}
        \end{center}
      \end{tarjeta}
    \end{column}
  \end{columns}
\end{frame}


% ── 3. LOS 5 COMPONENTES ─────────────────────────────────────────────────────
\begin{frame}{Los \alert{5 componentes}}
  \vspace{0.2cm}
  \begin{columns}[T]
    \begin{column}{0.44\textwidth}
      \begin{tarjeta}
        \textcolor{camarillo}{\faCircle[regular]}\enspace
        {\small\bfseries Solución parcial}\\
        {\footnotesize\color{cgris} Lista de estudiantes ya elegidos}
      \end{tarjeta}\vspace{4pt}
      \begin{tarjeta}
        \textcolor{camarillo}{\faCircle[regular]}\enspace
        {\small\bfseries Candidatos}\\
        {\footnotesize\color{cgris} Estudiantes del pool con índice mayor}
      \end{tarjeta}\vspace{4pt}
      \begin{tarjeta}
        \textcolor{camarillo}{\faCircle[regular]}\enspace
        {\small\bfseries \texttt{es\_valido()}}\\
        {\footnotesize\color{cgris} No supera el máx. por ciudad}
      \end{tarjeta}
    \end{column}
    \begin{column}{0.52\textwidth}
      \begin{tarjeta}
        \textcolor{camarillo}{\faCircle[regular]}\enspace
        {\small\bfseries \texttt{puede\_podar()}}\\
        {\footnotesize\color{cgris} Los candidatos restantes no alcanzan para cumplir un mínimo}
      \end{tarjeta}\vspace{4pt}
      \begin{tarjeta}
        \textcolor{camarillo}{\faCircle[regular]}\enspace
        {\small\bfseries Solución completa}\\
        {\footnotesize\color{cgris} Tamaño exacto + todos los mínimos cumplidos}
      \end{tarjeta}
    \end{column}
  \end{columns}
\end{frame}


% ── 4. LAS CUOTAS ─────────────────────────────────────────────────────────────
\begin{frame}{Las \alert{cuotas} a respetar}
  \vspace{0.3cm}
  \begin{columns}[c]
    \begin{column}{0.23\textwidth}
      \begin{tarjeta}
        \begin{center}
          {\normalsize\bfseries $\geq 2$}\\[2pt]
          {\footnotesize\color{cgris} de\\Computación}
        \end{center}
      \end{tarjeta}
    \end{column}
    \begin{column}{0.23\textwidth}
      \begin{tarjeta}
        \begin{center}
          {\normalsize\bfseries $\geq 1$}\\[2pt]
          {\footnotesize\color{cgris} de\\Matemáticas}
        \end{center}
      \end{tarjeta}
    \end{column}
    \begin{column}{0.23\textwidth}
      \begin{tarjeta}
        \begin{center}
          {\normalsize\bfseries $\leq 3$}\\[2pt]
          {\footnotesize\color{cgris} por\\ciudad}
        \end{center}
      \end{tarjeta}
    \end{column}
    \begin{column}{0.23\textwidth}
      \begin{tarjeta}
        \begin{center}
          {\normalsize\bfseries $\geq 2$}\\[2pt]
          {\footnotesize\color{cgris} de\\semestre 1}
        \end{center}
      \end{tarjeta}
    \end{column}
  \end{columns}
  \vspace{0.5cm}
  \begin{center}
    {\footnotesize\color{cgris} Conjunto 1 — Base}
  \end{center}
\end{frame}


% ── 5. LA PODA CLAVE ─────────────────────────────────────────────────────────
{
\setbeamercolor{background canvas}{bg=cverde}
\setbeamercolor{frametitle}{bg=cverde, fg=white}
\begin{frame}{La poda que lo cambió todo}
  \begin{center}
    {\normalsize Sin la \textbf{poda global de ciudades}:}\\[0.3cm]
    {\large el Conjunto 3 tardaba \alert{\textbf{8,380 llamadas}}}\\[0.1cm]
    {\large Con ella, termina en \alert{\textbf{1 sola llamada}}}
  \end{center}
  \vspace{0.5cm}
  \begin{formula}
    {\small\ttfamily\color{white}
    capacidad = $\sum$ (máx/ciudad $-$ actuales/ciudad)\\
    si capacidad $<$ faltantes $\;\rightarrow\;$ \textbf{imposible}}
  \end{formula}
\end{frame}
}


% ── 6. EL CÓDIGO ─────────────────────────────────────────────────────────────
\begin{frame}[fragile]{El algoritmo en \alert{Python}}
  \vspace{0.2cm}
  \begin{tcolorbox}[colback=ccode, colframe=white!10!ccode, arc=8pt, boxrule=0.5pt,
      left=8pt, right=8pt, top=6pt, bottom=6pt]
    \begin{lstlisting}[style=py]
def backtracking(pool, muestra, cuotas, soluciones, contadores, inicio=0):
    contadores["llamadas"] += 1
    if es_solucion_completa(muestra, cuotas):   # caso base
        soluciones.append(list(muestra))
        return
    if puede_podar(muestra, pool[inicio:], cuotas):  # poda
        return
    for i, candidato in enumerate(pool[inicio:]):
        if es_valido(muestra, candidato, cuotas):
            muestra.append(candidato)            # ELEGIR
            backtracking(pool, muestra, cuotas, soluciones, contadores, inicio+i+1)
            muestra.pop()                        # RETROCEDER
    \end{lstlisting}
  \end{tcolorbox}
\end{frame}


% ── 7. TABLA DE RESULTADOS ────────────────────────────────────────────────────
{
\setbeamercolor{background canvas}{bg=clight}
\setbeamercolor{normal text}{fg=cdark}
\setbeamercolor{frametitle}{bg=clight, fg=cdark}
\begin{frame}{3 conjuntos de \textcolor{cverde}{cuotas}}
  \vspace{0.4cm}
  \begin{center}
    \begin{tabular}{lcccc}
      \toprule
      \rowcolor{cverde}
      \color{white}\textbf{Conjunto} &
      \color{white}\textbf{Tamaño} &
      \color{white}\textbf{Llamadas} &
      \color{white}\textbf{Podas} &
      \color{white}\textbf{Resultado}\\
      \midrule
      C1 — Base         & 6 & 7          & 0          & \textcolor{cverde}{\bfseries\faCheckCircle\ Sí}\\
      C2 — Diversidad   & 6 & 7          & 0          & \textcolor{cverde}{\bfseries\faCheckCircle\ Sí}\\
      C3 — Sin solución & 7 & \textbf{1} & \textbf{1} & \textcolor{red!70!black}{\bfseries\faTimesCircle\ No}\\
      \bottomrule
    \end{tabular}
  \end{center}
\end{frame}
}


% ── 8. EXTENSIÓN ─────────────────────────────────────────────────────────────
\begin{frame}{La muestra \alert{más equilibrada}}
  \vspace{0.3cm}
  \begin{columns}[c]
    \begin{column}{0.44\textwidth}
      \begin{center}
        {\fontsize{42}{46}\selectfont\bfseries\color{camarillo} 10{,}734}\\[4pt]
        {\small\color{cgris} muestras válidas encontradas}
      \end{center}
    \end{column}
    \begin{column}{0.44\textwidth}
      \begin{center}
        {\fontsize{42}{46}\selectfont\bfseries\color{camarillo} 0.77}\\[4pt]
        {\small\color{cgris} score de balance máximo}
      \end{center}
    \end{column}
  \end{columns}
  \vspace{0.7cm}
  \begin{center}
    {\small\color{cgris}
      La mejor muestra incluye los \textbf{4 programas}:\\
      Comp.\,(2) · Mat.\,(2) · Est.\,(1) · Fís.\,(1)}
  \end{center}
\end{frame}


% ── 9. DECISIONES DEL GRUPO ──────────────────────────────────────────────────
\begin{frame}{Decisiones \alert{del grupo}}
  \vspace{0.2cm}
  \begin{columns}[T]
    \begin{column}{0.48\textwidth}
      \begin{tarjeta}
        {\small\bfseries Representación}\\
        {\footnotesize\color{cgris} Diccionarios Python — más legibles que índices numéricos}
      \end{tarjeta}\vspace{4pt}
      \begin{tarjeta}
        {\small\bfseries Bug encontrado}\\
        {\footnotesize\color{cgris} C3 tardaba 8,380 llamadas sin poda global de ciudades}
      \end{tarjeta}
    \end{column}
    \begin{column}{0.48\textwidth}
      \begin{tarjeta}
        {\small\bfseries Solución al bug}\\
        {\footnotesize\color{cgris} Condición 0: capacidad total de ciudades $<$ faltantes}
      \end{tarjeta}\vspace{4pt}
      \begin{tarjeta}
        {\small\bfseries Índice creciente}\\
        {\footnotesize\color{cgris} Evita combinaciones duplicadas sin memoria extra}
      \end{tarjeta}
    \end{column}
  \end{columns}
\end{frame}


% ── 10. EQUIPO ────────────────────────────────────────────────────────────────
\begin{frame}{¿Quién hizo \alert{qué?}}
  \vspace{0.3cm}
  \begin{columns}[c]
    \begin{column}{0.3\textwidth}
      \begin{tarjeta}
        \begin{center}
          {\large\faCode\par}\vspace{4pt}
          {\small\bfseries Sebastián\\Gutiérrez}\\[3pt]
          {\footnotesize\color{cgris} Algoritmo y poda}
        \end{center}
      \end{tarjeta}
    \end{column}
    \begin{column}{0.3\textwidth}
      \begin{tarjeta}
        \begin{center}
          {\large\faChartBar\par}\vspace{4pt}
          {\small\bfseries Mateo\\Bernal}\\[3pt]
          {\footnotesize\color{cgris} Dataset, gráficas y extensión}
        \end{center}
      \end{tarjeta}
    \end{column}
    \begin{column}{0.3\textwidth}
      \begin{tarjeta}
        \begin{center}
          {\large\faFileAlt\par}\vspace{4pt}
          {\small\bfseries Thomas\\Ramírez}\\[3pt]
          {\footnotesize\color{cgris} Informe, presentación y demo}
        \end{center}
      \end{tarjeta}
    \end{column}
  \end{columns}
\end{frame}


% ── 11. CONCLUSIONES ─────────────────────────────────────────────────────────
{
\setbeamercolor{background canvas}{bg=cverde}
\setbeamercolor{frametitle}{bg=cverde, fg=white}
\begin{frame}{Lo que aprendimos}
  \vspace{0.3cm}
  \begin{columns}[T]
    \begin{column}{0.3\textwidth}
      \begin{tarjetaverde}
        \begin{center}
          {\small\bfseries La poda es el motor}\\[4pt]
          {\footnotesize Sin ella: miles de llamadas.\\Con ella: 7.}
        \end{center}
      \end{tarjetaverde}
    \end{column}
    \begin{column}{0.3\textwidth}
      \begin{tarjetaverde}
        \begin{center}
          {\small\bfseries Imposible en 1 llamada}\\[4pt]
          {\footnotesize La poda global detecta sin explorar ninguna rama.}
        \end{center}
      \end{tarjetaverde}
    \end{column}
    \begin{column}{0.3\textwidth}
      \begin{tarjetaverde}
        \begin{center}
          {\small\bfseries Válida $\neq$ Óptima}\\[4pt]
          {\footnotesize 10,734 muestras válidas, solo una con balance 0.77.}
        \end{center}
      \end{tarjetaverde}
    \end{column}
  \end{columns}
\end{frame}
}


% ── 12. CIERRE ────────────────────────────────────────────────────────────────
\begin{frame}
  \begin{center}
    \vspace{1.2cm}
    {\large Un buen algoritmo no solo prueba caminos.}\\[0.4cm]
    {\Large\bfseries\color{camarillo} Aprende cuándo dejar de probarlos.}
    \vspace{1.5cm}

    {\footnotesize\color{cgris}
      UNAL Manizales · Intro a Ciencias de la Computación · 17 junio 2026}
  \end{center}
\end{frame}


\end{document}
